\section{The essential values criterion}
In this section, we~recall and state the standard notion of essential value, which is a classical tool to study skew products' ergodicity. Let $(X,\mathcal B, \mu)$ be a standard probability space. Let $T:X\to X$ be $\mu$-measure preserving automorphism and let $f:X\to \R^m$, where $m \geq 1$. Consider the skew product $T_f$ on $X\times \R^n$ given by
\[
T_f(x,r)=(Tx,r+f(x)).
\]
We say that $a\in\R^m$ is an \emph{essential value} of $T_f$ if for every $\epsilon>0$ and every measurable subset $E\subseteq X$ with $\mu(E)>0$, there exists $n\in\N$ such that
\[
\mu\{x\in E\mid T^{n}(x)\in E\text{ and }|S_n f(x)-a|<\epsilon\}>0,
\]
where $S_n f (x)$ denotes the $n$-th Birkhoff sum of $f$ evaluated at $x$, which is given by
\begin{equation}
\label{eq:Birkhoff_sums}
S_n f (x):=\begin{cases}
\sum_{i=0}^{n-1}f(T^{i}(x))& \text{ if }n\ge 1,\\
0&\text{ if }n=0,\\
-\sum_{i=-n}^{-1}f(T^{i}(x))&\text{ if }n\le 1.
\end{cases}
\end{equation}
We denote the set of essential values of $T_f$ by $\Ess(T_f)$.

The following classical fact links this notion to the ergodicity of the skew product.
\begin{theorem}\label{thm: essvalues}
With the notation as above, the set $\Ess(T_f)$ is a closed subgroup of~$\R^m$. Moreover, $T_f$ is ergodic with respect to $\mu\otimes \Leb_{\R^m}$ if and only if $T$ is ergodic and $\Ess(T_f)=\R^m$.
\end{theorem}
We will use a simplified version of this criterion, as introduced by Conze and Frączek in
\cite{conze_cocycles_2011}.

\begin{theorem}[{\cite[Lem.\,2.7]{conze_cocycles_2011}}]\label{thm: FUcriterion}
Let $a\in\R^m$. Assume that for every $\epsilon>0$ there exists a sequence
of subsets $\{\Xi_n\}_{n\in\N}$ and an increasing sequence
$\{q_n\}_{n\in\N}$ of natural numbers such that
\begin{enumerate}
\item \label{cond:measure} $\liminf_{n\to\infty}\mu(\Xi_n)>0$,
\item \label{cond:rigidity}$\lim_{n\to\infty}\sup_{x\in \Xi_n}|T^{q_n}(x)-x|=0$,
\item \label{cond:almost_invariant} $\lim_{n\to\infty}\mu(\Xi_n\triangle T(\Xi_n))=0$,
\item \label{cond:BS} $|S_{q_n}f(x)-a|<\epsilon$.
\end{enumerate}
Then $a\in \Ess(T_f)$.
\end{theorem}

Lemma 2.7 in \cite{conze_cocycles_2011} actually states that
the
topological support of the limit distribution
$P:=\lim_{n\to\infty}\frac{1}{\mu(\Xi_n)}(S_{q_n}f(x)|_{\Xi_n})_*\mu|_{\Xi_n}$,
which exists up
to taking a subsequence due to tightness guaranteed by Condition
\eqref{cond:BS}, is contained in $\Ess(T_f)$. However, again by \eqref{cond:BS},
the topological support of $P$ is contained in $\overline{B(a,\epsilon)}\subset \R^m$. By
passing with $\epsilon$ to 0 and by the fact that $\Ess(T_f)$ is a closed subset
of $\R^m$, we~get that $a\in \Ess(T_f)$.

We now provide a version of the above criterion, that is going to be effective for our purposes.

\begin{proposition}\label{prop: effcriterion}
Let $T:I\to I$ be an ergodic IET and let $m\in \N_+$. Assume that the function $f:I^{\times m}\to \R^{m}$ satisfies the following:
\begin{enumeratei}
\item \label{cond:continuity} $f$ is of the form $(x_1,\ldots,x_m)\mto (f_1(x_1),\ldots,f_m(x_m))$ with each $f_j$ being differentiable over the exchanged intervals,
\item \label{cond:sequences} there exist sequences $\{\Xi_n\}_{n\in\N}$ and $\{q_n\}_{n\in\N}$ satisfying \eqref{cond:measure}-\eqref{cond:almost_invariant} in Theorem \ref{thm: FUcriterion}, with the additional assumption that, for any $n \in \N$, $\Xi_n=\bigsqcup_{i=0}^{h_n-1}T^{i}({I_n})$ is a Rokhlin tower of $h_n\le q_n$ intervals with $|I_n|\le\sfrac{1}{q_n}$ and $\sup_{x\in \Xi_n}|T^{q_n}(x)-x|\le D/q_n$ for some $D>1$,
\item \label{cond:cont_interval} for every $x\in\Xi_n$, the interval $[x,T^{q_n}(x)]$ (or $[T^{q_n}(x),x]$) is a continuity interval of $f_j$, for every $j\in\{1,\ldots,m\}$,
\item \label{cond:derivative} there exists $C>1$ such that
\[
|S_{q_n}f_j(x)|\le C, \qquad C^{-1}q_n\le S_{q_n}f_j'(x)\le Cq_n,
\]
for any $j=1,\ldots,m$ and any $x\in\Xi_n$, and
\[
|f_j'(x)|\le C,
\]
for any $x\in I$.
\end{enumeratei}
If additionally $T^{\otimes m}$ is ergodic, then so is $T_f^{\otimes m}$.
\end{proposition}

\begin{proof}
In view of Theorem \ref{thm: essvalues}, it~is enough to show that we can
obtain an $m$\nobreakdash-dimensional cube of essential values. We will prove the
proposition for $m=1$ by realizing the essential values through Rokhlin
towers, which are subsets of $\Xi_n$. For $m\ge 2$ we first obtain the
same result for every $j=1,\ldots,m$, \ie we find an interval $(a_j,b_j)$
of essential values, realized through Rokhlin towers inside $\Xi_n$. Hence,
the set $(a_1,b_1)\times\ldots\times(a_m,b_m)$ is the set of essential
values realized via sequences of Rokhlin towers inside
$\bigsqcup_{i=0}^{h_n-1}(T^{\otimes m})^i(I_n^{\times m})$.

That being said, we~assume from now on that $m=1$. Note that due to our assumption on $\Xi_n$ and the derivative of $f$, each of the sets $S_{q_n}f(\Xi_n)\subseteq [-2C,2C]$ is a uniformly bounded union of intervals. Note that, since
\[
\liminf_{n\to\infty} h_n |I_n| = \liminf_{n\to\infty}\mu(\Xi_n) > 0
\]
and $h_n\le q_n$, we~have that, by~passing to a subsequence if necessary, there exists $E>1$ with
\[
E^{-1}\le q_n|I_n|\le 1.
\]
Hence, given the assumption on the derivative, we~have
\[
\Leb_{\R}\left(S_{q_n}f(\Xi_n)\right)\ge \frac{1}{CE}>0
\]
for every $n\in\N$.
By passing to a subsequence if necessary, we~may assume that the sets $S_{q_n}f(\Xi_n)$ have a nonempty intersection. Let $y\in \bigcap_{i=1}^{\infty} S_{q_n}f(\Xi_n)$. We claim that~$y$ is an essential value. For this purpose, we~will construct a subtower $\Xi_{n}^{y}\subseteq \Xi_n$, i.e., a~subset of $\Xi_n$ which itself is a Rokhlin tower of intervals, depending on $\epsilon>0$, such that \eqref{cond:measure} and \eqref{cond:BS} in Theorem \ref{thm: FUcriterion} is satisfied. This is enough, since \eqref{cond:rigidity} and \eqref{cond:almost_invariant} are automatically satisfied by any sequence of subtowers of $\Xi_n$.

Fix $\epsilon>0$ and consider the point {$x_n \in \Xi_n$ such that $S_{q_n}f(x_n)=y$}.
Let $\ell_n\in\{0,\ldots,h_n-1\}$ be such that $x_n\in T^{\ell_n}(I_n)$ and assume without loss of generality that $\ell_n<h_{n}/2$, the other case being treated symmetrically. Since each level of the tower is an interval then
\begin{align*}
\tag*{either}
\Bigl(x_n,x_n+\frac{1}{\max\{C, D,E\} q_n}\Bigr)\subseteq T^{\ell_n}(I_n),\\
\tag*{or} \Bigl(x_n-\frac{1}{\max\{C,D,E\} q_n},x_n\Bigr)\subseteq T^{\ell_n}(I_n).
\end{align*}
Again, without loss of generality, let us assume that it is the former. Consider the tower
\[
\Xi_n^y:=\bigsqcup_{i=0}^{\epsilon h_n/4CD} T^i\left(x_n,x_n+\frac{\epsilon}{2\max\{C,D,E\} q_n}\right)\subseteq \Xi_n.
\]
We now show that for every $x\in \Xi_n^y$ we have $S_{q_n}f(x)\in(y-\epsilon,y+\epsilon)$. If $x\in T^{\ell_n}(I_n)$, then by the mean value theorem, we~have
\[
\left|S_{q_n}f(x)-y\right|=\left|S_{q_n}f(x)-S_{q_n}f(x_n)\right|\le Cq_n|x-x_n|\le \epsilon/2.
\]
If $x\in T^{j} (T^{\ell_n}(I_n))$ with $j=1,\ldots,\epsilon h_n/4CD$, then we
split the Birkhoff sum into two pieces $S_{q_n}f(x) = S_{q_n - j}f (x) +
S_jf(T^{q_n - j}(x))$, which we estimate separately. For the first term, we~have
\[
\left|S_{q_n-j} f(x)-S_{q_n-j} f(T^j(x_n))\right|\le Cq_n|x-T^j(x_n)|\le\epsilon/2.
\]

For every $k\in\{0,\ldots,h_n/2\}$ we have that $T^{k}(x_n)$ and
$T^{k}(T^{-j}(x))$ belong to the same level of $\Xi_n$ and, since $\Xi_n$ is a
tower of intervals as such, belong to the same interval continuity of $f$.
Moreover, by~\eqref{cond:cont_interval}, $T^k(T^{q_n-j}(x))=T^k(T^{q_n}(T^{-j}(x)))$ and
$T^{k}(T^{-j}(x))$ belong to the same
continuity interval of $f$ for every $k\in\{0,\ldots,h_n/2\}$. Hence
$T^k(T^{q_n-j}(x))$ and $T^k(x_n)$ belong to the same continuity intervals of
$f$. Hence, by~\eqref{cond:BS} and mean value theorem we
get
\[
\left|S_j f(T^{q_n-j}(x))-S_j f(x_n)\right|\le \frac{C\epsilon h_n}{4D}\,|T^{q_n-j}(x)-x_n|\le \frac{C\epsilon h_n}{4CD} \frac{2D}{q_n}\le\epsilon/2
\]
and thus $|S_{q_n} f(x)-y|\le \epsilon$. It~remains to notice that
\begin{align*}
\liminf_{n\to\infty}\Leb(\Xi_n^y)&\ge \liminf_{n\to\infty}\frac{\epsilon h_n}{4CD}\frac{1}{\max\{C,D,E\}}\,|I_n|\\
&\ge \liminf_{n\to\infty} \frac{1}{4(\max\{C,D,E\})^2}\Leb(\Xi_n)>0.
\end{align*}
Thus we have proved that $y\in \Ess(T_f)$. Note that for every $n\in\N$ {we have}
\[
x\in \Bigl(x_n,x_n+\frac{1}{2\max\{C,D,E\} q_n}\Bigr).
\]
Hence,
\[
S_{q_n}f(x)-S_{q_n}f(x_n)\ge C^{-1}q_n |x-x_n|.
\]
Therefore, by~applying similar reasoning as to $y$, we~obtain that every
\[
z\in\Bigl[y,y+\frac{1}{2C^{-1}\max\{C,D,E\}}\Bigr]
\]
is an essential value of $T_f$. This finishes the proof of the proposition.
\end{proof}
\section{Proofs of main results.}

This section contains the proof of Theorems \ref{thm: ergo}, \ref{thm: uniqergo} and \ref{thm: index}. In all three proofs, we~will apply the ergodicity criterion described in Proposition \ref{prop: effcriterion}. For this reason, we~start this section by outlining a construction that will be common to all of the proofs since it concerns only the underlying IET $T$ and not the cocycle being considered. At the end of this construction, we~will describe in detail why the assumptions of Proposition \ref{prop: effcriterion} are fulfilled in each setting.

Throughout this section, let $T=(\pi,\lambda): I \to I$ be an ergodic symmetric IET on $d = \#\A$ intervals $\{I_\alpha\}_{\alpha \in \A}$.

By Corollary \ref{cor:centers_connections}, there exists $\beta \in \A$ such that $c_{\beta}$ is not a part of any connection of~$T$. By Lemma \ref{lem:sym_int_disjoint}, there exists $\alpha \in \A$ and a nested sequence of $\beta$-symmetric intervals $\{J_n\}_{n \in \N}$ disjoint from the connections of $T$, with endpoints $T^{-m_n}(\partial I_{\alpha})$ and $T^{m_n}(\partial I_{\hat\alpha})$ for some $m_n \nearrow \infty$, where $\pi_0(\hat\alpha) = \pi_0(\alpha)-1$, and such that $\{c_\beta\} = \bigcap_{n \in \N} J_n$.

By Proposition \ref{prop:induced_symmetric}, for every $n\in\N$, the induced IET $T_{J_n}$ is a symmetric IET with $d - d'$ intervals, where $d'$ is the number of non-trivial irreducible connections of~$T$. Without loss of generality we may index their exchanged intervals using the same alphabet $\B \subseteq \A$. Let us denote by $\{I_\gamma^n\}_{\gamma \in \B}$ the intervals exchanged by $T_{J_n}$ and by $\{c_\gamma^n\}_{\gamma \in \B}$ their middle points, where, to avoid the use of double subscripts, we~changed our usual notation $I_\gamma^{J_n}$ (\resp $c_\gamma^{J_n}$) to $I_\gamma^n$ (\resp $c_\gamma^n$).

By Proposition \ref{prop:induced_centers}, each of the towers in the tower decomposition of $I$ associated with $T_{J_n}$
\begin{equation}
\label{eq:decomposition_sequence}
I = \bigsqcup_{\gamma \in \B} \bigsqcup_{i = 0}^{h_\gamma^n - 1} T^i(I_\gamma^n),
\end{equation}
where $h^n = (h^n_\gamma)_{\gamma \in \B}$ is some vector in $\N^\A_+$ (see Section \ref{sc:induced_IETs}), contains exactly one point from the set
\begin{multline*}
 \{c _\sigma \mid \sigma \in \A \cup \{\sfrac{1}{2}\},\, \sigma \neq \beta \text{ and } c_\sigma \text{ does not belong to}\\
\text{any non-trivial irreducible connection}\},
\end{multline*}
in the middle of its \emph{central level} $T^{\lfloor h_\gamma^n/2
\rfloor}(I^n_\gamma)$. More precisely, for every $\gamma \in \B$ there exists
$\sigma$ in the set above such that the middle point $c_\gamma^n$ of the
interval $I_\gamma^n$ verifies $c_\sigma = T^{\lfloor h_\gamma^n/2
\rfloor}(c_\gamma^n)$.

Using these facts, we~will show how to build towers $\{\Xi_n\}_{n\in\N}$ and a sequence $\{q_n\}_{n\in\N}$ that satisfies the assumptions of Proposition \ref{prop: effcriterion}.

A natural approach would be to consider subtowers of the already constructed towers. However in their current form, the towers may be very unbalanced: the wide towers may be very short and thus of very small measure, while thin towers may be very tall and contain the majority of the interval $I$. Since we need to construct towers of measure bounded away from 0, we~would have to choose them to be inside of the thin towers. This however makes it very difficult to control the rigidity of $T$ inside those towers as well as to estimate the values of Birkhoff sums. We tackle this problem by jumping between the points around which we induce.

Consider a Rokhlin tower $X_n:=\bigsqcup_{i=0}^{h_{\gamma_n}^n-1}T^i(I^n_{\gamma_n})$, where $\gamma_n \in\mathcal \B$ is chosen so that its Lebesgue measure is the largest compared to the other towers in the decomposition \eqref{eq:decomposition_sequence}. In particular
\[
\Leb(X_n)\geq \frac{1}{\#\B}\ge\frac{1}{\#\A}.
\]

Let us denote by $\mathfrak J_n$ the central level of this tower and recall that it contains a point $c_\sigma$ for some $\sigma \in \A \cup \{\sfrac{1}{2}\}$. By Proposition \ref{prop:induced_symmetric}, the induced transformation $T_{\mathfrak J_n}$ is a symmetric IET with $d - d'$ intervals, and we denote its exchanged intervals by $\{\mathfrak I_{\gamma}^n\}_{\gamma \in \mathcal B}$. As before, we~have a decomposition in Rokhlin towers of the form
\[
 I = \bigsqcup_{\gamma \in \B} \bigsqcup_{i = 0}^{\mathfrak h_\gamma^n - 1} T^i(\mathfrak I_\gamma^n).
\]

Let $\Gamma_n \in \B$ be such that $\mathfrak I_{\Gamma_n}^n$ is the largest of all intervals exchanged by $T_{\mathfrak J_n}$. Up to taking a subsequence, let us assume without loss of generality that there exists $\Gamma \in \B$ such that $\Gamma_n = \Gamma$, for every $n \in \N$.

As before, by~Proposition \ref{prop:induced_centers}, the tower $\bigsqcup_{i=0}^{\mathfrak h_{\Gamma}^n-1}T^i(\mathfrak I_{\Gamma}^n)$ contains exactly one point of the form $c_{\sigma}$ for some $\sigma \in \A \cup \{\sfrac{1}{2}\}$ in the middle of its central level. Let us denote this point by $\mathfrak c_n$.

Define
\[
\Xi_n:=\bigsqcup_{i=0}^{h_{\gamma_n}^n/2-1}T^{i}(\mathfrak I_{\Gamma}^n)\qquad \text{and}\qquad q_n:=\mathfrak h_{\Gamma}^n.
\]

Before passing to the proofs of each of the theorems, let us show that $\{\Xi_n\}_{n\in\N}$ and $\{q_n\}_{n\in\N}$ satisfy the assumptions \eqref{cond:sequences} and \eqref{cond:cont_interval} in Proposition \ref{prop: effcriterion}.

We start by showing that \eqref{cond:sequences} in Proposition \ref{prop: effcriterion} is satisfied, that is, that the sequences above verify \eqref{cond:measure}--\eqref{cond:almost_invariant} in Theorem \ref{thm: FUcriterion}.

First, we~can easily check that $\{\Xi_n\}_{n \in \N}$ satisfies \eqref{cond:measure} in Theorem \ref{thm: FUcriterion}. Indeed,
\begin{align*}
\Leb(\Xi_n) & =(h_{\gamma_n}^n/2)|\mathfrak I_{\Gamma}^n| \\ & > \frac{1}{\#\A}(h_{\gamma_n}^n/2)|\mathfrak J_n|=\frac{1}{\#\A}(h_{\gamma_n}^n/2)|I^n_{\gamma}| \\
& > \frac{1}{3\#\A} h_{\gamma_n}^n|I^n_{\gamma}|=\frac{1}{3\#\A} \Leb(X_n) \\
& >\frac{1}{3\#\A^2}.
\end{align*}

To see that $\{\Xi_n\}_{n\in\N}$ satisfies \eqref{cond:almost_invariant} in Theorem \ref{thm: FUcriterion} it is enough to notice that
\[
\mu(\Xi_n\triangle T(\Xi_n))\le 2|\mathfrak I_{\Gamma}^n|\to 0 \quad\text{as }n\to\infty.
\]

We will now show that
\begin{equation}\label{eq: derbound}
\sup_{x\in\Xi_n}|T^{q_n}(x)-x|<D/q_n\quad\text{for some }D>1
\end{equation}
thus showing \eqref{cond:rigidity} in Theorem \ref{thm: FUcriterion} as well. The argument uses the same observation as the one used to prove \eqref{cond:cont_interval} in Proposition \ref{prop: effcriterion}, which is the following.

Let $j\in\{0,\ldots,h_{\gamma_n}^n/2\}$ and let $x\in T^j(\mathfrak I_{\Gamma}^n)$. Then $T^{q_n-j}(x)\in \mathfrak J_n$. However, $\mathfrak J_n$ is the middle level of the tower $X_n$. Hence we have that $x$
and $T^{q_n}(x)=T^{j}(T^{q_n-j}(x))$ belong to the same level of the tower $X_n$. In particular, they belong to a continuity interval of $T$ (and as such to a continuity interval of any function continuous over exchanged intervals). Moreover, we~have
\[
|T^{q_n}(x)-x|\le |\mathfrak J_n|\le \#\A|\mathfrak I_{\Gamma}^n|=
\#\A\frac{1}{q_n} q_n|\mathfrak I_{\Gamma}^n|\le \frac{\#\A}{q_n}.
\]
Thus Conditions \eqref{cond:sequences} and \eqref{cond:cont_interval} in Proposition \ref{prop: effcriterion} are satisfied.

In the following proofs, we~will check that the remaining assumptions in Proposition~\ref{prop: effcriterion}, namely, Conditions \eqref{cond:continuity} and \eqref{cond:derivative}, are satisfied for the different cocycles considered in each setting. In view of the construction above this is enough to apply Proposition \ref{prop: effcriterion} and conclude the ergodicity of the skew product under consideration.

\begin{proof}[Proof of Theorem \ref{thm: ergo}]
We assume that $a>0$ since the other case follows symmetrically. We will apply Proposition \ref{prop: effcriterion} for $m=1$. Since $f(x)=a(x-\sfrac{1}{2})$ is continuous, \eqref{cond:continuity} in \ref{prop: effcriterion} is satisfied. We now show that \eqref{cond:derivative} is satisfied.

First, trivially, we~have
\[
f'(x)=a\quad\text{ for }x\in I.
\]
In particular
\[
S_{q_n} f'(x)=aq_n\quad \text{ for }x\in I.
\]
Recall that the tower $\bigsqcup_{i=0}^{\mathfrak h_{\Gamma}^n-1}T^i(\mathfrak I_{\Gamma}^n)$ has $\mathfrak c_n$ as its central point. By construction, the first visit time of $\mathfrak c_n$ via $T^{-1}$ to $\mathfrak J_n$ is $q_{n}/2$ or $(q_n+1)/2$. In both cases, by~Lemma \ref{lem: berktrujillo},
\[
S_{q_n}f(T^{-\lfloor (q_n+1)/2\rfloor}(\mathfrak c_n))=0.
\]
Since $S_{q_n}f$ is continuous in $\mathfrak I_{\Gamma}^n$ and $q_n|\mathfrak I_{\Gamma}^n|<1$, by~the mean value theorem we have that
\[
\left|S_{q_n}f(x)\right|<a\quad \text{ for every }x\in \mathfrak I_{\Gamma}^n.
\]
If $x\in T^j(\mathfrak I_{\Gamma}^n)$ for $j\in \{1,\ldots,h_{\gamma_n}^n/2\}$, then by \eqref{eq: derbound} and, again, by~the mean value theorem,
\[
\left|S_{q_n}f(x)-S_{q_n}f(T^{-j}(x)) \right|=
\left|S_j f(T^{q_n-j}(x))-S_j f(T^{-j}(x)) \right|
\le a D.
\]
Thus we get
\[
\left|S_{q_n}f(x)\right|<a(D+1) \quad \text{ for every }x\in\Xi_n.
\]
Since the bound does not depend on $n$, \eqref{cond:derivative} is satisfied and, by~Proposition \ref{prop: effcriterion}, the skew product $T_f$ is ergodic.
\end{proof}
